\section*{Capítulo \ref{ch:b1:mammal-organization} --- Organización funcional de un mamífero}

\begin{solution}{exo:b1:mammal-organization:1}
De vertebrado: simetría bilateral, \omterm{def:b1:mammal-organization:bodyplan}{cefalización}, eje vertebral dorsal que
encierra la médula espinal, \omterm{def:b1:mammal-organization:bodyplan}{celoma} ventral con las vísceras, cuatro
extremidades. De mamífero: el \omterm{def:b1:mammal-organization:bodyplan}{diafragma} que divide el \omterm{def:b1:mammal-organization:bodyplan}{celoma} en tórax y
abdomen, el pelo, la leche, la temperatura corporal regulada.
\end{solution}

\begin{solution}{exo:b1:mammal-organization:2}
Nutrición: digestivo (estómago), respiratorio (pulmón), circulatorio
(corazón), excretor (riñón). Relación: nervioso (encéfalo), \omterm{def:b1:mammal-organization:organ}{órganos} de los
sentidos (ojo), endocrino (tiroides), esqueleto y músculos (fémur, bíceps),
piel. Reproducción: genital (ovario, testículo), glándulas mamarias.
\end{solution}

\begin{solution}{exo:b1:mammal-organization:3}
Agua $0.6\times 60 = \qty{36}{L}$; intracelular \qty{24}{L};
extracelular \qty{12}{L}, de los cuales intersticial \qty{9}{L} y \omterm{def:b1:mammal-organization:compartments}{plasma}
\qty{3}{L}.
\end{solution}

\begin{solution}{exo:b1:mammal-organization:4}
$1.1/5 = 22\%$. \qty{1.1}{L/min} son \qty{1580}{L} al día, es decir, el
volumen sanguíneo 290 veces.
\end{solution}

\begin{solution}{exo:b1:mammal-organization:5}
Inulina restante \qty{1.7}{g}; $V = 1.7/0.12 = \qty{14}{L}$ de
extracelular; intracelular $42 - 14 = \qty{28}{L}$.
\end{solution}

\begin{solution}{exo:b1:mammal-organization:6}
La pared capilar deja pasar libremente el agua y los iones pequeños, que se
equilibran, pero retiene las proteínas. Las proteínas plasmáticas atraen
agua hacia los vasos por ósmosis; sin ellas el agua saldría al espacio
intersticial, que se hincharía (edema), y el volumen plasmático
disminuiría.
\end{solution}

\begin{solution}{exo:b1:mammal-organization:7}
En paralelo, cada \omterm{def:b1:mammal-organization:organ}{órgano} recibe sangre con todo el oxígeno y toda la presión
arteriales, y su flujo puede regularse de forma independiente. En serie, el
encéfalo recibiría sangre ya empobrecida por los músculos, a menor presión,
y su aporte caería precisamente durante el ejercicio.
\end{solution}

\begin{solution}{exo:b1:mammal-organization:8}
$hS\Delta T = 80 - 20 = \qty{60}{W}$; $h = 60/(1.8\times 13) =
\qty{2.6}{W.m^{-2}.K^{-1}}$. Con el aire a \qty{33}{\celsius} la pérdida no
evaporativa se anula; la piel debe calentarse por encima del aire (la
dilatación vascular la sube hacia \qty{36}{\celsius}) y la sudoración debe
llevarse el resto.
\end{solution}

\begin{solution}{exo:b1:mammal-organization:9}
Calor a evacuar $800 - 200 - 150 = \qty{450}{W}$, es decir, \qty{1620}{kJ/h},
\qty{675}{g} de sudor por hora. Sin sudoración, $450/(70\times 3500) =
\qty{1.8e-3}{K/s}$, alrededor de \qty{0.11}{K/min}: \qty{3}{K} en media hora.
\end{solution}

\begin{solution}{exo:b1:mammal-organization:10}
Los cuerpos pequeños pierden calor más deprisa por gramo (gran $S/V$) y
necesitan una producción continua, alta y fiable: la grasa parda produce
calor químicamente a ritmo constante y sin movimiento, mientras que la
tiritona exige músculos desarrollados (de los que carecen los recién
nacidos), estorba el movimiento y la alimentación, y es intermitente. La
grasa parda se sitúa además junto a los grandes vasos y calienta
directamente la sangre.
\end{solution}

\begin{solution}{exo:b1:mammal-organization:11}
El hipotálamo calentado lee ``demasiado calor'': el perro jadea y dilata sus
vasos cutáneos pese al frío, y su temperatura central cae uno o dos grados
en una hora. Sin hipotálamo no hay \omterm{def:b1:organism-environment:homeostasis}{valor de consigna}: el perro ni tirita con
el frío ni jadea con el calor, y su temperatura deriva hacia la de la
habitación en ambos casos.
\end{solution}

\begin{solution}{exo:b1:mammal-organization:12}
El \omterm{def:b1:mammal-organization:compartments}{líquido extracelular} (\omterm{def:b1:mammal-organization:compartments}{líquido intersticial} y \omterm{def:b1:mammal-organization:compartments}{plasma}), de composición rica
en sodio, temperatura y volumen casi constantes, baña cada célula como el
agua de mar bañaba a los primeros animales. Su constancia la mantienen
bucles de \omterm{def:b1:organism-environment:homeostasis}{retroalimentación negativa} (el riñón para el volumen y los iones,
el pulmón para los gases, el hígado y el páncreas para la glucosa, el
hipotálamo para la temperatura). Sus ``orillas'' son las \omterm{prop:b1:organism-environment:surfaces}{superficies de
intercambio} --- tubo digestivo, pulmón, riñón, piel ---, donde se encuentra
con el exterior, y la circulación es la corriente que lo agita.
\end{solution}

\begin{solution}{pb:b1:mammal-organization:1}
\textbf{1.} $B = 3.4\times 2^{0.75} = \qty{5.7}{W}$.
\textbf{2.} Madriguera: $5.7\times 12\times 3600 = \qty{246}{kJ}$;
búsqueda de alimento: $17.1\times 6\times 3600 = \qty{369}{kJ}$; al aire
libre: $11.4\times 6\times 3600 = \qty{246}{kJ}$.
\textbf{3.} \qty{861}{kJ}, es decir, $861/493 = 1.75$ veces la basal
($B$ a lo largo de \qty{24}{h} $= \qty{493}{kJ}$).
\textbf{4.} $0.25\times 18\times 0.65 = \qty{2.9}{kJ/g}$.
\textbf{5.} $861/2.9 = \qty{295}{g}$ de hierba fresca, \qty{74}{g} de
materia seca.
\textbf{6.} $295/2000 = 15\%$ de su masa.
\textbf{7.} $861/20 = \qty{43}{L}$ de $\mathrm{O_2}$; \qty{30}{mL/min}.
\textbf{8.} $0.9\times 43 = \qty{39}{L}$ de $\mathrm{CO_2}$, es decir, $39/22.4\times 44 = \qty{76}{g}$.
\textbf{9.} $0.75\times 295 = \qty{221}{g}$.
\textbf{10.} Materia seca digerida $0.65\times 74 = \qty{48}{g}$; agua
$0.6\times 48 = \qty{29}{g}$.
\textbf{11.} No digerida \qty{26}{g}; heces $26/0.6 = \qty{43}{g}$, que
se llevan \qty{17}{g} de agua.
\textbf{12.} $(44 - 5)\times 0.86 = \qty{34}{g}$.
\textbf{13.} Entradas: $221 + 29 = \qty{250}{g}$. Salidas: $17 + 130 + 34 + 20
= \qty{201}{g}$. Excedente \qty{49}{g}.
\textbf{14.} No: la hierba aporta más de lo que pierde; el excedente sale
como orina adicional (unos \qty{180}{mL} en total).
\textbf{15.} La misma materia seca, \qty{74}{g}, en $\qty{148}{g}$ de
hierba: agua del alimento \qty{74}{g}; entradas \qty{103}{g} frente a
salidas \qty{201}{g}: un déficit de \qty{100}{g} al día, que debe beber o
perder de su cuerpo.
\textbf{16.} Agua $0.6\times 2.0 = \qty{1.2}{L}$; extracelular
\qty{0.4}{L}; \omterm{def:b1:mammal-organization:compartments}{plasma} \qty{0.1}{L}.
\textbf{17.} $1.5/15 = \qty{0.10}{L}$: concuerda.
\textbf{18.} $0.10/0.60 = \qty{0.167}{L}$, el \qty{8}{\%} de la masa corporal.
\textbf{19.} $5\times(2/70)^{0.75} = \qty{0.35}{L/min}$; tiempo de
circuito $0.167/0.35 = \qty{0.48}{min}$, unos \qty{29}{s}.
\textbf{20.} $350/200 = \qty{1.75}{mL}$ por latido.
\textbf{21.} $0.1\times 2^{2/3} = \qty{0.16}{m^2}$.
\textbf{22.} $2.5\times 0.16\times 34 = \qty{13.6}{W}$, frente a los
$2B = \qty{11.4}{W}$ supuestos: el conejo debe producir algo más, o
reducir la pérdida con la postura y el refugio.
\textbf{23.} Abiertas: $10\times 0.01\times 34 = \qty{3.4}{W}$, otra
cuarta parte de la pérdida por el pelaje; cerradas: $10\times 0.01\times 3 =
\qty{0.3}{W}$. Las cierra.
\textbf{24.} Pelaje: $2.5\times 0.16\times 9 = \qty{3.6}{W}$; orejas
abiertas: $10\times 0.01\times 9 = \qty{0.9}{W}$; total \qty{4.5}{W}
frente a los \qty{5.7}{W} producidos. Los \qty{1.2}{W} restantes deben
salir por evaporación: el conejo jadea (\qty{1.2}{W} son \qty{1.8}{g} de
agua por hora), se tumba sobre el suelo fresco y se queda a la sombra.
\textbf{25.} Gasto energético diario \qty{0.86}{MJ}, $1.75$ veces su tasa
basal; recambio de agua \qty{250}{g}, el \qty{12.5}{\%} de su masa.
\end{solution}
